\documentclass[a4paper,11pt]{article}
\usepackage[utf8]{inputenc}
\usepackage[T1]{fontenc}
\usepackage[english]{babel}
\usepackage[left=2cm,right=2cm,top=2cm,bottom=2cm]{geometry}
\usepackage{xcolor}
\usepackage{hyperref}
\usepackage{fontawesome5}
\usepackage{parskip}

% --- Couleurs ---
\definecolor{primary}{RGB}{0, 82, 147}
\hypersetup{colorlinks=true, urlcolor=primary}

\begin{document}

% --- EXPÉDITEUR ---
\noindent
\begin{minipage}[t]{0.5\textwidth}
    \textbf{Loïc Aron MBASSI EWOLO}\\
    Engineering Student -- ENSEA / Polytechnique Yaoundé\\[3pt]
    \faMapMarker\ Cergy, France\\
    \faPhone\ (+33) 7 59 52 33 98\\
    \faEnvelope\ \href{mailto:loicaron.mbassiewolo@ensea.fr}{loicaron.mbassiewolo@ensea.fr}
\end{minipage}
\hfill
% --- DATE ---
\begin{minipage}[t]{0.4\textwidth}
    \raggedleft
    Cergy, February 5th, 2026
\end{minipage}

\vspace{1.2cm}

% --- DESTINATAIRE ---
\hfill
\begin{minipage}[t]{0.5\textwidth}
    \raggedleft
    \textbf{Mines Saint-Étienne}\\
    Flexible Electronics Department\\
    Prof. Marc RAMUZ\\
    Gardanne, France
\end{minipage}

\vspace{1cm}

\noindent\textbf{Subject:} Application for Internship -- Wireless Implanted Sensor for Animals' Sleep Study

\vspace{0.8cm}

Dear Professor Ramuz,

Building instrumented hardware that interfaces with biological signals has been a driving interest throughout my studies. When I read about your project on implantable wireless sensors and optical energy transfer through skin, I recognized a unique opportunity to apply my electronics and embedded systems background to biomedical research.

At the IoT Laboratory of Polytechnique Yaoundé, I participated in the development of instrumented gloves for sign language recognition. This project involved sensor integration (IMU, flex sensors), firmware development in C, component soldering, and hardware assembly into functional enclosures. I also gained experience with test bench automation and rigorous technical documentation. During my UROP project mentored by Google DeepMind researchers, I worked on ECG signal analysis for cardiac anomaly detection, which gave me a solid foundation in biomedical signal processing and scientific methodology.

Currently pursuing a dual degree at ENSEA in embedded systems and electronics, I am developing skills in signal processing, simulation (MATLAB/Simulink), and hardware design. My coursework covers the fundamentals needed for your project: electronic circuit design, sensor interfaces, and experimental validation. I am comfortable with laboratory work and eager to learn optical simulation tools like COMSOL.

I am particularly drawn to the interdisciplinary nature of your research, combining flexible electronics, optical energy transfer, and biomedical applications. The opportunity to work in your cleanroom facilities and contribute to a multi-partner project on implantable sensors would be an invaluable experience for my future career in bioelectronics.

I am available from April 2026 for 4 to 6 months and would welcome the opportunity to discuss how I can contribute to your team.

Thank you for considering my application.

\vspace{0.8cm}

Sincerely,

\vspace{0.5cm}

\textbf{Loïc Aron MBASSI EWOLO}\\
\textit{Engineering Student -- ENSEA / Polytechnique Yaoundé}

\end{document}
